\chapter{Installation and updates}

Installers are on the releases page:
\web{github.com/nobrega8/modbench/releases}. Each release lists what changed.

\section{Windows}

Download \file{Modbench\_x.y.z\_x64-setup.exe} (recommended) or the
\file{.msi}, and run it. The installer is about 2~MB: Modbench uses the
WebView2 that comes with Windows.

\section{macOS}

Download \file{Modbench\_x.y.z\_universal.dmg}: one app for Apple Silicon and
Intel Macs. Open it and drag Modbench to \nav{Applications}.

\begin{note}[First start]
Modbench is not signed with an Apple Developer ID, so macOS asks before the
first start. Right-click the app in \nav{Applications}, choose \nav{Open}, then
\nav{Open} again. Later starts need nothing.
\end{note}

\section{Linux}

\begin{itemize}
  \item \textbf{AppImage}: download it, make it executable
        (\val{chmod +x Modbench*.AppImage}) and run it.
  \item \textbf{Debian, Ubuntu and derivatives}: install the \file{.deb}, or
        add the APT repository so that \val{apt upgrade} brings new versions.
        It has packages for \val{amd64} (PCs), \val{arm64} (Raspberry Pi 3,
        4 and 5 with a 64-bit system, and other 64-bit ARM boards) and
        \val{armhf} (32-bit Raspberry Pi OS), for Debian 12 and 13, Ubuntu
        22.04 and later and Raspberry Pi OS Bookworm or later; APT picks the
        one for the machine. Raspberry Pi OS Buster and older cannot run
        Modbench (it lacks WebKitGTK 4.1).
\end{itemize}

\begin{lstlisting}[language=bash]
curl -fsSL https://modbench.nobrega.me/key.asc | sudo gpg --dearmor -o /usr/share/keyrings/modbench.gpg
echo "deb [signed-by=/usr/share/keyrings/modbench.gpg] https://modbench.nobrega.me stable main" | sudo tee /etc/apt/sources.list.d/modbench.list
sudo apt update && sudo apt install modbench
\end{lstlisting}

\begin{note}[Serial ports on Linux]
To open serial ports without root, add your user to the \val{dialout} group
(\val{sudo usermod -aG dialout \$USER}) and log in again.
\end{note}

\section{Updates}

Modbench checks for a new version when it starts (this can be turned off in
\nav{Settings\then Updates}). When there is one, a banner says what changed:

\screenshot[2.2cm]{update-banner}{The update banner.}
  {The banner at the top of the window: \emph{Modbench x.y.z is available},
   a summary of the changes, \emph{Install and restart} and \emph{Later}.}

\uiaccent{Install and restart} downloads the update, checks its signature,
installs it and starts Modbench again. \ui{Later} hides the banner until the
next start. \nav{Settings\then Updates} can also check by hand.

Installs from the APT repository update with the system
(\val{sudo apt upgrade}) instead.

\section{Files Modbench opens}

The installers register two file types, each with its own icon. A
double-click opens them in Modbench, also when it is already running.

\begin{center}
\begin{tabularx}{\linewidth}{l L}
  \toprule
  \thead{File} & \thead{What it is} \\
  \midrule
  \file{.mbmap.json} & A device map (Chapter~\ref{ch:maps}). Opens on the Main page. \\
  \file{.mbmap} & A device map in the older version 1 format; Modbench can convert it. \\
  \file{.mbsim} & A simulator setup (Section~\ref{sec:sim-setups}). Opens in the simulator. \\
  \bottomrule
\end{tabularx}
\end{center}

\begin{note}
Windows tells file types apart by their last extension only, so
\file{.mbmap.json} files keep the icon of JSON files there. On Linux they get
the map icon.
\end{note}
